\section*{Chapter \ref{ch:s2:delta-one-and-dividends} --- Delta-One and Dividends}

\begin{solution}{exo:s2:delta-one-and-dividends:1}
$(7 - 5)$ basis points of \$50 million: \$10\,000.
\end{solution}

\begin{solution}{exo:s2:delta-one-and-dividends:2}
$(40 - 15) \times 0.25 - 5 = 1.25$ basis points of \$100 million: \$12\,500 for the quarter.
\end{solution}

\begin{solution}{exo:s2:delta-one-and-dividends:3}
$108.7 / 95.7 - 1 = 13.6\%$ in total, $1.136^{1/3} - 1 = 4.3\%$ a year. The simulated mean, 3.8\% a year, is lower because recession years cut
dividends and the average of annualised returns is below the annualised average.
\end{solution}

\begin{solution}{exo:s2:delta-one-and-dividends:4}
The note's buyer is short a put on the index through the protection; the issuer is long it and hedges by holding the index. Higher dividends lower the
index's forward and raise the put's value to the issuer, lowering the note it owes: a 0.1 point rise in the dividend yield gains it 0.11 per 100.
\end{solution}

\begin{solution}{exo:s2:delta-one-and-dividends:5}
If issuers of structured products sell short-dated dividends to hedge, the buyers who take them are paid a discount, which shows up as a high expected
return on short-term dividend claims; a risk premium and a supply premium are hard to tell apart in prices alone.
\end{solution}

\begin{solution}{exo:s2:delta-one-and-dividends:6}
Banks report and are charged on their balance sheets at quarter-ends, so they shrink them then; the few that can hold the basket against the future
demand more for doing so, and the future's implied financing rate rises.
\end{solution}

\begin{solution}{exo:s2:delta-one-and-dividends:7}
Mean P\&L per trade is 1.14 basis points at 1 minute, 0.77 at 2, 0.40 at 3, 0.05 at 4 and $-0.26$ at 5: the trade stops paying between four and five
minutes. With a half-life of 10 minutes, each minute of delay gives back part of a gap that is only a few basis points beyond the costs.
\end{solution}

\begin{solution}{exo:s2:delta-one-and-dividends:8}
The return comes with equity risk: the one-year trade correlates 0.41 with the index, loses in about a fifth of years and 24.4\% on average in a
recession year, when bonds usually do well. It is paid for supplying insurance to structured-product issuers, not for lending.
\end{solution}

\begin{solution}{pb:s2:delta-one-and-dividends:1}
\begin{enumerate}
\item Trading an index future against its basket when it departs from fair value by more than the cost of both, unwinding when the gap closes.
\item Minute mispricings with a 4 basis point standard deviation and 10-minute half-life; 2 basis points a side for the basket and 0.5 for the
future, 5 for a round trip.
\item 2\,230 trades a year: 1.64 basis points each without lag, 1.14 at one minute, 0.77 at two, $-0.26$ at five.
\item It enters beyond the round-trip cost and leaves at zero, so each trade earns at least the excess.
\item Holding an asset against a derivative that delivers it later, earning the implied financing spread less funding and capital.
\item The price above spot less dividends implies the rate at which the market finances the basket.
\item At the roll: 10\% of quarters, 15.4 basis points locked, 0.4 a year; in the quarter-end window: 67.5\%, 40.7, 8.0 a year.
\item The bank holds the basket and pays the client its return for a spread over funding.
\item The note's buyer is short a put; the issuer's hedge makes it long dividends: 0.11 per 100 for 0.1 point of yield, about 1.1 million per billion.
\item The discount of dividend futures to expected dividends when natural holders of dividends sell more than others will buy at fair value.
\item EUR 100 per point of the EURO STOXX 50's dividends in the contract year, cash settled, up to ten years out.
\item Futures from 96.7 to 93.7 against expected dividends from 100.8 to 117.2; mean returns 3.8--4.2\% a year; the one-year trade's standard deviation
12.3\%, 5\% quantile $-25.5\%$, chance of a loss 20.9\%.
\item Held to expiry, the one-year dividend future earns 4.2\% a year on average from a 4\% supply discount, with a loss in about one year in five
and $-24.4\%$ in a recession year.
\item From 2004 dividend swaps were a major profit source for hedge funds because implied dividends traded at high discounts and dividends grew.
\item Short-term dividend claims have higher expected returns, Sharpe ratios and volatilities than the index, with betas below one.
\item A two-factor model fits the dividend futures curve, and investors revise dividends beyond the business cycle only a little.
\item In recessions, when dividends are cut and equities fall together.
\item Total-return swap financing.
\item Both earn a spread for providing balance sheet, and both widen when banks' balance sheets are scarce.
\item Balance sheet and hedging capacity: the index's return without funding it, and dividends' risk without the index's.
\end{enumerate}
\end{solution}

\begin{solution}{iq:s2:delta-one-and-dividends:1}
Spot less the present value of dividends to expiry, grown at the financing rate to expiry. Dividend forecasts, the financing rate the market uses
(not the risk-free rate), withholding tax and the timing of ex-dates can all be wrong.
\end{solution}

\begin{solution}{iq:s2:delta-one-and-dividends:2}
Bottom up: each constituent's announced and forecast dividends, payout policies and earnings, weighted by index weights and ex-dates; checked
against dividend futures and option-implied forwards.
\end{solution}

\begin{solution}{iq:s2:delta-one-and-dividends:3}
From the cost of holding the basket over the period: the bank's funding plus the balance-sheet charge, which is higher over quarter-end, plus
dividends passed through and the cost of trading the basket; compare with the futures-implied financing over the same dates.
\end{solution}

\begin{solution}{iq:s2:delta-one-and-dividends:4}
A recession with dividend cuts of 30--50\% in the near years, correlated with an equity fall, and a regulatory ban on bank dividends; margin calls on
the futures as prices fall.
\end{solution}

\begin{solution}{iq:s2:delta-one-and-dividends:5}
A feed of index weights, corporate actions and dividend forecasts; financing curves; real-time spot of the basket; recomputation on each input change;
alerts when the future departs from fair value; and versioned inputs for audit.
\end{solution}

\begin{solution}{iq:s2:delta-one-and-dividends:6}
$F(D + \delta) - F(D) = -\delta e^{rT}$. A trader who underestimates dividends by $\delta$ computes a fair value too high by $\delta e^{rT}$, sees the
future as cheap by that amount and buys it: the apparent arbitrage is the forecast error, and it loses when the dividends are paid.
\end{solution}
